\section{到達点：正規な実際の退化因子を持つ場合の分類}

基礎体は $\CC$ とする。引き続き、$\PGL_5$ による商を取る前の全射影 Jacobi スキーム
\[
 \cP=\{[\pi]\in\PP H^0(\PP^4,\wedge^2T_{\PP^4}):[\pi,\pi]=0\}
\]
を対象とする。本稿では一般階数四の点を考え、
\[
 D=V(\pi\wedge\pi),\qquad Y=\operatorname{Sing}(D),\qquad C=Y_{\mathrm{red}}
\]
と置く。$D$ は実際の退化五次因子である。
以下の正規性は $D$ 自体に対する条件であり、被約化 $D_{\mathrm{red}}$ だけに課す条件ではない。
正規性の下で $D$ は被約かつ既約である。実際、異なる超曲面成分は射影空間内で交わり、
その交点の局所環は整域でない。

\Needspace{12\baselineskip}
\begin{theorem}[正規な退化因子の場合：複素閉点の分類]\label{thm:normal}
第20・21版で証明した帰着を入力とする。
一般階数四の $\PP^4$ 上の Poisson 構造 $\pi$ の実際の退化因子 $D$ が正規ならば、
$C$ は滑らかな楕円正規五次曲線であり、
\[
 D=\operatorname{Sec}(C).
\]
射影座標変換と非零定数倍を除き、$\pi$ はこの楕円曲線に付随する
Feigin--Odesskii の $q_{5,1}$ 型である。
\end{theorem}
逆に、滑らかな楕円正規五次曲線から得られる $q_{5,1}$ はこの条件を満たす。
その退化因子は割線五次超曲面で、特異支持は楕円曲線であり、正規である\cite{Pym}。

この定理の新しい部分は、第21版に残った八つの $\widetilde E_8$ 支持を全て排除することである。
既知の purely elliptic な場合の分類\cite{Pym}と、本研究の追加の排除を区別する。
Pym の定理自体が一般の正規退化因子を分類している、と主張するものではない。
また定理\ref{thm:normal}は複素閉点の分類であり、Jacobi スキームの非被約構造の決定ではない。

\subsection{前版から引き継ぐ厳密な入力}
第20・21版の全ファイルと PDF を変更せずに同梱した。入力は次の通りである。
\begin{enumerate}
\item $Y$ は純一次元 Cohen--Macaulay で $\deg Y=40$。
$C$ は連結、$\pi|_C=0$ であり、標準二次持ち上げの Pfaffian が $F$ の五つの偏微分を与える。
\item modular residue は $C$ の各既約成分で恒等的に零ではない。
全 $\widetilde E_6$ 型の楕円五次の場合は $q_{5,1}$ 型として完結している。
それ以外では全成分の一般横断型が $\widetilde E_8$、$\deg C=4$ である。
\item 四次の場合、各成分の正規化は $\PP^1$。
modular residue は支持の特異点の全ての逆像で零となり、
各成分の標点は高々二つである。
\item 四次の場合、超平面に含まれる全ての支持と、既約な支持は排除済みである。
残る $C$ は $\PP^4$ を張る、算術種数零の被約可約四次曲線である。
その特異点は独立な接線を持つ滑らかな枝の普通の多重点であり、
incidence graph は木になる。
\item 残る支持の射影同値類は以下の八つで尽きる。
各配置は単一の射影軌道であり、曲線の配置について未処理の連続パラメータはない。
\end{enumerate}
これらの入力を今回の計算で再証明したとは述べない。
前版の支持の網羅性と今回の各支持の排除とを合わせて、定理\ref{thm:normal}を得る。

\subsection{八つの配置と今回の排除}
$e_i$ を座標点とし、
\[
 L_{ij}=\langle e_i,e_j\rangle,\quad
 A_{ijk}=[s^2e_i+st e_j+t^2e_k],\quad
 T_{0123}=[s^3e_0+s^2te_1+st^2e_2+t^3e_3]
\]
と書く。配置を表す数字は成分の次数の分割であり、Poisson スキームの成分番号ではない。

\begin{center}\small
\begin{tabular}{@{}p{23mm}p{64mm}p{47mm}@{}}\toprule
型 & 被約特異支持 $C$ & 排除の主な理由\\\midrule
$1111$ 鎖 & $L_{01}\cup L_{12}\cup L_{23}\cup L_{34}$ & 二直線の交点は二重 modular 零点\\
$1111$ 三重点と枝 & $L_{01}\cup L_{02}\cup L_{03}\cup L_{14}$ & 同上、中央直線の零因子次数\\
$1111$ 四重点 & $L_{01}\cup L_{02}\cup L_{03}\cup L_{04}$ & 全 slope が零、五次式が錐になる\\
$211$ conic 中央 & $A_{012}\cup L_{03}\cup L_{24}$ & conic 上に二つの二重零点\\
$211$ 直線中央 & $A_{012}\cup L_{03}\cup L_{34}$ & 中央直線の零因子次数\\
$211$ 三重点 & $A_{012}\cup L_{03}\cup L_{04}$ & 三つの係数の短い線形結合\\
$22$ 一節点 & $A_{012}\cup A_{034}$ & 平面 Hessian と全例外値の係数証明書\\
$31$ 外部直線 & $T_{0123}\cup L_{04}$ & residue の全場合と非正規な二例外族\\\bottomrule
\end{tabular}\end{center}

\subsection{今回の結論の範囲}
全 $\PP^4$ の Poisson 構造の分類には、一般階数二、非正規な被約退化因子、
非被約退化因子の各場合が残る。
また族の閉包、既約成分同士の交わり、Jacobi スキームの冪零構造も今回の定理の外である。
前版までの全体の既約成分数の下界57は変更せず、正確な個数とはしない。
今回残る二つの非正規な Jacobi 解族も、それだけで新しい既約成分とは数えない。
